%! AP Statistics — Unit 4, Lesson 4.11 — Homework
\documentclass[10pt]{article}
\usepackage{apstats-article}
\usepackage{apstats-boxes}

\begin{document}

\pageheader{Unit 4, Lesson 4.11}{Homework}
\namedateperiod

\vspace{0.1in}

\begin{practicebox}
\small
\textbf{Directions:} For each problem, show all formula setup before computing.
Write a complete-sentence interpretation wherever asked.

\vspace{10pt}
\textbf{1.} A seed company claims its tomato seeds have an 85\% germination rate.
A gardener plants 30 seeds. Let $X =$ the number that germinate.
Assume $X \sim B(30, 0.85)$.
\begin{enumerate}[label=(\alph*), leftmargin=2em, itemsep=8pt]
  \item Calculate $\mu_X$ and $\sigma_X$. Show formula setup.\\
        \writeline\\[1pt]\writeline
  \item Interpret $\mu_X$ in context.\\
        \writeline\\[1pt]\writeline
  \item Find $P(X \le 22)$. Write the calculator notation.\\
        \writeline
  \item Find $P(X \ge 28)$. Write the inequality rewrite and the calculator notation.\\
        \writeline\\[1pt]\writeline
\end{enumerate}

\vspace{8pt}
\textbf{2.} A multiple-choice test has 50 questions, each with 5 options.
A student guesses on every question. Let $X =$ the number correct.
\begin{enumerate}[label=(\alph*), leftmargin=2em, itemsep=8pt]
  \item State the distribution of $X$: \blank{4cm}
  \item Compute $\mu_X$ and $\sigma_X$.\\
        \writeline
  \item The passing score is 30 out of 50. Find $P(X \ge 30)$.
        Is this likely to happen by guessing?\\
        \writeline\\[1pt]\writeline
\end{enumerate}

\vspace{8pt}
\textbf{3.} A hospital reports that 40\% of emergency room patients are discharged
within 4 hours. On a given day, 120 patients visit the ER.
Let $X =$ the number discharged within 4 hours. Assume $X \sim B(120, 0.40)$.
\begin{enumerate}[label=(\alph*), leftmargin=2em, itemsep=8pt]
  \item $\mu_X = $ \blank{2cm} \qquad $\sigma_X \approx$ \blank{2cm}
  \item Interpret $\mu_X$: \writeline\\[1pt]\writeline
  \item Find $P(40 \le X \le 50)$. Write this as a difference of two \texttt{binomcdf} values.\\
        \writeline\\[1pt]\writeline
\end{enumerate}

\vspace{8pt}
\textbf{4.} (AP-Style Free Response) \textit{According to a survey, 62\% of adults in
a large city support a new public transit plan. A researcher randomly samples
45 adults (less than 10\% of the city's adult population).
Let $X =$ the number who support the plan.}
\begin{enumerate}[label=(\alph*), leftmargin=2em, itemsep=8pt]
  \item Define the distribution of $X$ and check all BINS conditions briefly.\\
        \writeline\\[1pt]\writeline\\[1pt]\writeline
  \item Calculate and interpret $\mu_X$.\\
        \writeline\\[1pt]\writeline
  \item Calculate $\sigma_X$.\\
        \writeline
  \item Find $P(X \ge 30)$. Show the equality rewrite and the calculator expression.\\
        \writeline\\[1pt]\writeline
\end{enumerate}

\vspace{8pt}
\textbf{5. Extension:} Use technology or a table to compute $P(X \ge 1)$ for
$X \sim B(20, 0.05)$ two ways: (i) as $1 - P(X = 0)$ and (ii) using
\texttt{binomcdf}. Verify the answers agree and explain why these two approaches
are equivalent.

\vspace{4pt}
\writeline\\[1pt]\writeline\\[1pt]\writeline

\end{practicebox}

\end{document}
